% ============================================================
\chapter*{Pr\'eface}
\addcontentsline{toc}{chapter}{Pr\'eface}
% ============================================================

L'apprentissage automatique (\emph{machine learning}) est une discipline
\`a l'intersection des math\'ematiques, de l'informatique et de la
statistique. Son objectif fondamental est de concevoir des algorithmes
capables d'\textbf{apprendre \`a partir de donn\'ees}~: au lieu de
programmer explicitement une r\`egle de d\'ecision, on laisse l'algorithme
d\'ecouvrir cette r\`egle \`a partir d'exemples.

\section*{Positionnement du cours}

Ce cours se situe au niveau L3/Master~1 et suppose des pr\'erequis en~:
\begin{itemize}
  \item \textbf{Alg\`ebre lin\'eaire}~: espaces vectoriels, matrices,
    valeurs propres, d\'ecomposition en valeurs singuli\`eres.
  \item \textbf{Probabilit\'es et statistiques}~: variables al\'eatoires,
    esp\'erance, variance, loi normale, estimation, tests.
  \item \textbf{Analyse}~: d\'eriv\'ees partielles, gradient, optimisation
    de fonctions de plusieurs variables.
  \item \textbf{Programmation Python}~: NumPy, Matplotlib ; la connaissance
    de scikit-learn est un plus mais n'est pas requise.
\end{itemize}

\section*{Philosophie p\'edagogique}

Chaque chapitre suit une progression en trois temps~:
\begin{enumerate}
  \item \textbf{Th\'eorie}~: d\'erivation math\'ematique rigoureuse des
    algorithmes, avec \'enonc\'e et preuve des r\'esultats cl\'es.
  \item \textbf{Pratique}~: impl\'ementation en Python avec scikit-learn
    et, quand c'est instructif, \`a partir de z\'ero.
  \item \textbf{Analyse}~: discussion des forces, faiblesses, hypoth\`eses
    et cas d'utilisation de chaque m\'ethode.
\end{enumerate}

\section*{Plan du cours}

\begin{center}
\begin{tikzpicture}[
  block/.style={rectangle, draw, rounded corners, minimum width=4cm,
    minimum height=0.8cm, align=center, font=\small},
  arr/.style={-{Stealth[length=2.5mm]}, thick, steelblue}
]
  \node[block, fill=blue!10] (ch1) {Ch.~1~: Introduction};
  \node[block, fill=blue!10, below=0.3cm of ch1] (ch2) {Ch.~2~: R\'egression};
  \node[block, fill=blue!10, below=0.3cm of ch2] (ch3) {Ch.~3~: Classification};
  \node[block, fill=green!10, below=0.3cm of ch3] (ch4) {Ch.~4~: R\'egularisation};
  \node[block, fill=green!10, below=0.3cm of ch4] (ch5) {Ch.~5~: SVM};
  \node[block, fill=green!10, below=0.3cm of ch5] (ch6) {Ch.~6~: Arbres};

  \node[block, fill=orange!10, right=2cm of ch1] (ch7) {Ch.~7~: Ensembles};
  \node[block, fill=orange!10, below=0.3cm of ch7] (ch8) {Ch.~8~: Non supervis\'e};
  \node[block, fill=orange!10, below=0.3cm of ch8] (ch9) {Ch.~9~: R\'eduction dim.};
  \node[block, fill=red!10, below=0.3cm of ch9] (ch10) {Ch.~10~: Bay\'esien};
  \node[block, fill=red!10, below=0.3cm of ch10] (ch11) {Ch.~11~: S\'election};
  \node[block, fill=red!10, below=0.3cm of ch11] (ch12) {Ch.~12~: Noyaux};

  \draw[arr] (ch1) -- (ch2);
  \draw[arr] (ch2) -- (ch3);
  \draw[arr] (ch3) -- (ch4);
  \draw[arr] (ch4) -- (ch5);
  \draw[arr] (ch5) -- (ch6);
  \draw[arr] (ch6.east) -- ([xshift=0.5cm]ch6.east) |- (ch7.west);
  \draw[arr] (ch7) -- (ch8);
  \draw[arr] (ch8) -- (ch9);
  \draw[arr] (ch9) -- (ch10);
  \draw[arr] (ch10) -- (ch11);
  \draw[arr] (ch11) -- (ch12);
\end{tikzpicture}
\end{center}

\noindent Les chapitres 1--6 couvrent les fondements (supervis\'e lin\'eaire
et non lin\'eaire). Les chapitres 7--9 abordent les m\'ethodes d'ensemble
et l'apprentissage non supervis\'e. Les chapitres 10--12 traitent de sujets
avanc\'es (approche bay\'esienne, s\'election de mod\`eles, m\'ethodes \`a
noyaux).

\vfill
\noindent\textit{Bonne lecture et bon apprentissage~!}
